\chapter{Kod źródłowy i Potok Przetwarzania (Pipeline)}

Struktura kodu źródłowego została zaprojektowana w sposób modularny, co pozwala
na ścisłą separację logiki przetwarzania sygnałów, ekstrakcji cech oraz wizualizacji rezultatów.
Całość implementacji opiera się na spójnym potoku przetwarzania danych (ang. \textit{pipeline}), 
od surowych odczytów z sensorów, aż po gotowe wektory cech i ich analizę.

\section{Przepływ sterowania w projekcie}
Przepływ danych opiera się na liniowej sekwencji przetwarzania, w której wyjście
jednego etapu staje się wejściem kolejnego. Wizualizacja przepływu danych prezentuje się następująco:
% Zaktualizowano odniesienie do diagramu
\begin{figure}[H]
    \centering
    \includegraphics[width=0.8\linewidth]{flowchart.png}
    \caption{Architektura potoku przetwarzania danych (Pipeline)}
    \label{fig:pipeline_flowchart}
\end{figure}

Główny przepływ realizowany jest poprzez następujące po sobie etapy: 
\texttt{main} $\rightarrow$ \texttt{data\_processing} $\rightarrow$ \texttt{konwersja} $\rightarrow$ \texttt{segmentacja} $\rightarrow$ \texttt{ekstrakcja\_cech}.

Poniżej szczegółowo opisano każdy z etapów zgodnie z najlepszymi praktykami dokumentacji technicznej.

\section{Etap 1: Inicjalizacja i Sterowanie (\texttt{main})}
\textbf{Nazwa modułu/funkcji:} \texttt{main()} \\
\textbf{Ogólna idea:} Punkt wejściowy aplikacji. Odpowiada za orkiestrację całego procesu, inicjalizację połączeń (np. z bazą danych lub systemem plików) oraz sekwencyjne wywoływanie kolejnych etapów potoku przetwarzania. \\
\textbf{Wejście:} Parametry konfiguracyjne (ścieżki do archiwów ZIP, flagi sterujące, dane dostępowe do API). \\
\textbf{Wyjście:} Wywołanie procedur przetwarzających i przekazanie im uchwytów do danych (np. obiekty \texttt{MeasurementDatasetHook}).

\section{Etap 2: Przetwarzanie danych (\texttt{data\_processing})}
\textbf{Nazwa modułu/funkcji:} \texttt{process\_file()} / \texttt{data\_processing()} \\
\textbf{Ogólna idea:} Moduł odpowiedzialny za wczytanie surowych plików JSONL z archiwum, przeprowadzenie wstępnego czyszczenia (obsługa brakujących wartości) i przygotowanie struktury \texttt{ADC\_data}, która przechowuje znormalizowane wyniki pomiarów w woltach. Realizuje też podstawową filtrację (np. \texttt{outlier\_detection()}) w celu usunięcia szumów i anomalii pomiarowych. \\
\textbf{Wejście:} Surowe pliki pomiarowe w formacie JSONL pobrane z API lub archiwum lokalnego. \\
\textbf{Wyjście:} Ustrukturyzowany, ciągły sygnał czasowy oczyszczony z największych artefaktów.

\section{Etap 3: Konwersja sygnału (\texttt{konwersja})}
\textbf{Nazwa modułu/funkcji:} \texttt{convert\_signal()} / \texttt{breath\_separation()} \\
\textbf{Ogólna idea:} Analiza lokalnych ekstremów ciągłego sygnału w celu jego podziału na pojedyncze cykle oddechowe. Sygnał napięciowy jest rzutowany na fizyczne wartości oddychania, a faza wdechu i wydechu zostaje precyzyjnie oznaczona. \\
\textbf{Wejście:} Oczyszczony, ciągły sygnał czasowy (wynik etapu \texttt{data\_processing}). \\
\textbf{Wyjście:} Zbiór oznaczonych, pojedynczych cykli oddechowych (fale).

\section{Etap 4: Segmentacja (\texttt{segmentacja})}
\textbf{Nazwa modułu/funkcji:} \texttt{split\_data\_into\_segments()} \\
\textbf{Ogólna idea:} Dzielenie ciągłego sygnału na stałe, 2-minutowe okna (segmenty) umożliwiające analizę statystyczną i ujednolicenie danych wejściowych dla algorytmów klasyfikacji. Segmenty o utracie danych przekraczającej próg (np. 20\%) są odrzucane. Długość próbek wewnątrz segmentów jest normalizowana (resampling). \\
\textbf{Wejście:} Wyodrębnione cykle oddechowe. \\
\textbf{Wyjście:} Pakiety (segmenty) znormalizowanych danych gotowe do bezpośredniej ekstrakcji cech.

\section{Etap 5: Ekstrakcja cech (\texttt{ekstrakcja\_cech})}
\textbf{Nazwa modułu/funkcji:} \texttt{basic\_feature\_extraction()} / Skrypt \texttt{extract\_features} \\
\textbf{Ogólna idea:} Przekształcenie sygnału w dziedzinie czasu w niskowymiarowy wektor cech (np. BPM, głębokość oddechu, czas trwania fazy wdechu/wydechu). Funkcja jest wywoływana pojedynczo dla każdego zaakceptowanego segmentu sygnału. \\
\textbf{Wejście:} Znormalizowane segmenty sygnału (wynik etapu \texttt{segmentacja}). \\
\textbf{Wyjście:} Macierz cech (wektory cech) zapisana do pliku (np. \texttt{extracted\_features.jsonl}), stanowiąca bezpośrednie wejście dla modeli uczenia maszynowego.

\section{Moduły Analizy i Wizualizacji}

Poza głównym potokiem przetwarzania, system implementuje narzędzia badawcze i ewaluacyjne:

\subsection{Tworzenie profili oddechowych (\texttt{create\_data\_profiles})}
Odpowiada za tworzenie indywidualnych profili oddechowych (modele referencyjne). Zwraca mapę zawierającą najbardziej reprezentatywny sygnał dla danej osoby, z uwzględnieniem \texttt{timestamps}, \texttt{inhale}, \texttt{exhale}, \texttt{length} oraz \texttt{depth}. Wykorzystuje funkcję \texttt{get\_mode\_breath()} do wyłonienia dominanty oddechowej.

\subsection{Wizualizacja danych i redukcja wymiarowości (\texttt{visualize\_data})}
Moduł ładuje wyekstrahowane cechy (struktura \texttt{FeatureData}) i aplikuje na nich algorytmy redukcji wymiarowości:
\begin{itemize}
    \item \textbf{PCA (\texttt{PCA\_algorithm})} - analiza głównych składowych.
    \item \textbf{MDS (\texttt{MDS\_algorithm})} - wielowymiarowe skalowanie do celów wizualnych.
\end{itemize}
Moduł posiada także \texttt{SVM\_validation} służącą do wstępnej weryfikacji liniowej separowalności danych na płaszczyźnie, oraz metodę \texttt{plotheatmap} do wizualizacji cech po standaryzacji do przedziału $[0,1]$.

\section{Komunikacja z Infrastrukturą Zewnętrzną}

\subsection{Obsługa zapytań API (\texttt{DatabaseHandler})}
Klasa \texttt{DatabaseHandler} służy jako warstwa abstrakcji (Data Access Layer) przy wymianie informacji z BRICS API. Implementuje wzorzec CRUD:
\begin{itemize}
	\item \texttt{upload\_measurement()}: Zapis nowego pomiaru oraz metadanych do bazy w chmurze.
	\item \texttt{download\_measurement()}: Pobieranie zasobów po unikalnym identyfikatorze lub zestawie filtrów.
	\item \texttt{delete\_measurement()}: Przesłanie żądania usunięcia.
\end{itemize}

\subsection{Zarządzanie archiwalnymi zestawami danych (\texttt{MeasurementDatasetHook})}
Narzędzie systemowe do masowego procesowania danych. Realizuje interfejs iteratora w celu automatycznego wyodrębniania pomiarów z archiwów ZIP:
\begin{minted}{python}
mdh = MeasurementDatasetHook(target="raw")
for file in mdh:
    process_file(file)
\end{minted}
Dzięki takiemu podejściu aplikacja nie ładuje całości danych do pamięci RAM, zachowując wysoką wydajność.
